\section{Tool Generalization：MCP 时代的新训练难题}

讲者在后半段提出了一个现实挑战：用户可通过 MCP 或自定义 API 接入任意工具，训练时无法穷举全部工具集合。因此，模型要学习的不是固定工具调用模板，而是 \textbf{工具泛化能力（tool generalization）}。

\begin{knowledgebox}{工具泛化的三层能力}
\begin{enumerate}
    \item 读取并理解新工具描述（schema、参数、约束）；
    \item 在多工具竞争下选择正确工具；
    \item 调用失败后能自我修正并切换策略。
\end{enumerate}
\end{knowledgebox}

这与传统 function-calling benchmark 的差异在于，后者通常在封闭集合中评测；而真实产品里工具集合是开放、动态、个性化的。模型必须依赖强 instruction following 与反思能力才能稳健适配。

\begin{importantbox}{为什么 instruction following 是工具泛化前提}
讲者指出：如果模型连基本指令都难以稳定遵循，那么面对新工具描述时，策略会迅速退化。先修好 instruction following，再谈工具泛化，通常更有效。
\end{importantbox}

此外，课程中提到 ``未按提示使用指定工具要惩罚''。这体现了一个关键训练原则：工具调用正确性不只看结果，也看过程是否满足显式约束（例如必须使用某工具获取权威数据）。

\begin{warningbox}{仅靠工具调用成功率评估会误导}
模型可能通过偶然路径得到正确答案，但违反了产品合规流程。若评估缺少过程约束，部署后风险会显著上升。
\end{warningbox}

\subsection{本章小结}
在 MCP 时代，agent 模型的竞争力来自 ``对未知工具的迁移能力''。这要求训练同时强化 instruction following、反思纠错和过程合规。

